% =============================================================================
% Kapitel 2: Grundlagen
% -----------------------------------------------------------------------------
% Hier werden die Begriffe und Konzepte geklärt, die der Hauptteil voraussetzt.
% Beachte: Zwischen \section und der ersten \subsection stehen 2--4 Sätze
% (Kapitel-Einleitung), die den Inhalt zusammenfassen und an Kapitel 1 anknüpfen.
% =============================================================================
\section{Grundlagen}
\label{sec:grundlagen}

Dieses Kapitel klärt die zentralen Begriffe der Arbeit und ordnet sie in den
Forschungsstand ein. Es schafft damit die Grundlage für die Analyse in
Kapitel~\ref{sec:hauptteil}. Behandelt werden zunächst die begrifflichen
Grundlagen und anschließend der relevante Stand der Forschung.

% --- 2.1 ---
\subsection{Begriffliche Grundlagen}
\label{subsec:begriffe}

Definiere hier die Schlüsselbegriffe. Eine Definition wird belegt:
Ein \ac{KMU} ist nach gängiger Definition \dots\autocite[Vgl.][S.~3]{musterinstitut2023}

% --- 2.2 ---
\subsection{Stand der Forschung}
\label{subsec:forschungsstand}

Hier ordnest du die wichtigsten Quellen ein und zeigst, wo deine Arbeit
ansetzt. So bindest du eine Abbildung ein (Datei liegt unter abbildungen/):

% Beispiel fuer eine Abbildung -- erst einkommentieren, wenn ein Bild vorliegt,
% und im main.tex \listoffigures aktivieren.
% \begin{figure}[H]
%   \centering
%   \includegraphics[width=0.8\textwidth]{beispielbild.png}
%   \caption{Beschreibung der Abbildung}
%   \label{fig:beispiel}
% \end{figure}
% Im Text immer referenzieren: Abbildung~\ref{fig:beispiel} zeigt \dots

So bindest du eine Tabelle ein:

\begin{table}[H]
  \centering
  \caption{Beispielhafte Gegenüberstellung}
  \label{tab:beispiel}
  \begin{tabularx}{\textwidth}{@{}lX@{}}
    \toprule
    \textbf{Kriterium} & \textbf{Beschreibung} \\
    \midrule
    Merkmal A & Erläuterung zu Merkmal A \\
    Merkmal B & Erläuterung zu Merkmal B \\
    \bottomrule
  \end{tabularx}
\end{table}

\noindent Tabelle~\ref{tab:beispiel} stellt die Kriterien gegenüber.
